% Options for packages loaded elsewhere
\PassOptionsToPackage{unicode}{hyperref}
\PassOptionsToPackage{hyphens}{url}
\PassOptionsToPackage{dvipsnames,svgnames,x11names}{xcolor}
\documentclass[
  10pt,
  fontset=fandol]{ctexart}
\usepackage{xcolor}
\usepackage[margin=16mm]{geometry}
\usepackage{amsmath,amssymb}
\setcounter{secnumdepth}{-\maxdimen} % remove section numbering
\usepackage{iftex}
\ifPDFTeX
  \usepackage[T1]{fontenc}
  \usepackage[utf8]{inputenc}
  \usepackage{textcomp} % provide euro and other symbols
\else % if luatex or xetex
  \usepackage{unicode-math} % this also loads fontspec
  \defaultfontfeatures{Scale=MatchLowercase}
  \defaultfontfeatures[\rmfamily]{Ligatures=TeX,Scale=1}
\fi
\usepackage{lmodern}
\ifPDFTeX\else
  % xetex/luatex font selection
\fi
% Use upquote if available, for straight quotes in verbatim environments
\IfFileExists{upquote.sty}{\usepackage{upquote}}{}
\IfFileExists{microtype.sty}{% use microtype if available
  \usepackage[]{microtype}
  \UseMicrotypeSet[protrusion]{basicmath} % disable protrusion for tt fonts
}{}
\makeatletter
\@ifundefined{KOMAClassName}{% if non-KOMA class
  \IfFileExists{parskip.sty}{%
    \usepackage{parskip}
  }{% else
    \setlength{\parindent}{0pt}
    \setlength{\parskip}{6pt plus 2pt minus 1pt}}
}{% if KOMA class
  \KOMAoptions{parskip=half}}
\makeatother
\setlength{\emergencystretch}{3em} % prevent overfull lines
\providecommand{\tightlist}{%
  \setlength{\itemsep}{0pt}\setlength{\parskip}{0pt}}
\usepackage{amsmath,amssymb,mathtools,needspace}
\xeCJKDeclareCharClass{CJK}{"2032,"0400->"04FF,"2160->"216B,"2460->"2473,"25A0,"25C7,"25CB,"25CF}
\IfFontExistsTF{Arial Unicode MS}{\xeCJKsetup{AutoFallBack=true}\setCJKfallbackfamilyfont{\CJKrmdefault}{Arial Unicode MS}}{}
\setcounter{secnumdepth}{0}
\setlength{\emergencystretch}{2em}
\allowdisplaybreaks
\usepackage{bookmark}
\IfFileExists{xurl.sty}{\usepackage{xurl}}{} % add URL line breaks if available
\urlstyle{same}
\hypersetup{
  pdftitle={化高阶方程组为一阶方程组},
  colorlinks=true,
  linkcolor={Maroon},
  filecolor={Maroon},
  citecolor={Blue},
  urlcolor={Blue},
  pdfcreator={LaTeX via pandoc}}

\title{化高阶方程组为一阶方程组}
\author{}
\date{}

\begin{document}
\maketitle

\subsubsection{5.6.2 化高阶方程组为一阶方程组}\label{ux5316ux9ad8ux9636ux65b9ux7a0bux7ec4ux4e3aux4e00ux9636ux65b9ux7a0bux7ec4}

关于高阶微分方程（或方程组）的初值问题，原则上总可以归结为一阶方程组来求解。例如，考察下列 \(m\) 阶微分方程：

\[
y^{(m)}=f(x,y,y',\cdots,y^{(m-1)}),
\tag{5.6.3}
\]

初始条件为

\[
y(x_0)=y_0,\quad y'(x_0)=y'_0,\quad\cdots,\quad y^{(m-1)}(x_0)=y_0^{(m-1)}.
\tag{5.6.4}
\]

只要引进新的变量 \(y_1=y,y_2=y',\cdots,y_m=y^{(m-1)}\)，即可将 \(m\) 阶方程（5.6.3）化为如下的一阶方程组：

\[
\begin{cases}
y'_1=y_2,\\
y'_2=y_3,\\
\qquad\vdots\\
y'_{m-1}=y_m,\\
y'_m=f(x,y_1,y_2,\cdots,y_m);
\end{cases}
\tag{5.6.5}
\]

初始条件（5.6.4）则相应地化为

\href{../../source-images/pdf-147.jpeg}{原页}

\[
y_1(x_0)=y_0,\quad y_2(x_0)=y'_0,\quad\cdots,\quad y_m(x_0)=y_0^{(m-1)}.
\tag{5.6.6}
\]

不难证明，初值问题（5.6.3）、（5.6.4）和（5.6.5）、（5.6.6）是彼此等价的。

特别地，对于下列二阶方程的初值问题：

\[
\begin{cases}
y''=f(x,y,y'),\\
y(x_0)=y_0,\quad y'(x_0)=y'_0.
\end{cases}
\]

引进新的变量 \(z=y'\)，即可化为下列一阶方程组的初值问题：

\[
\begin{cases}
y'=z,\quad y(x_0)=y_0,\\
z'=f(x,y,z),\quad z(x_0)=y'_0.
\end{cases}
\]

针对这个问题应用四阶 Runge-Kutta 格式（5.6.1），有

\[
\begin{cases}
\displaystyle y_{n+1}=y_n+\frac h6(K_1+2K_2+2K_3+K_4),\\[6pt]
\displaystyle z_{n+1}=z_n+\frac h6(L_1+2L_2+2L_3+L_4).
\end{cases}
\]

按式（5.6.2），有

\[
K_1=z_n,\quad L_1=f(x_n,y_n,z_n);
\]

\[
K_2=z_n+\frac h2L_1,\quad
L_2=f\left(x_n+\frac h2,y_n+\frac h2K_1,z_n+\frac h2L_1\right);
\]

\[
K_3=z_n+\frac h2L_2,\quad
L_3=f\left(x_n+\frac h2,y_n+\frac h2K_2,z_n+\frac h2L_2\right);
\]

\[
K_4=z_n+hL_3,\quad L_4=f(x_n+h,y_n+hK_3,z_n+hL_3).
\]

如果消去 \(K_1,K_2,K_3,K_4\)，则上述格式可表示为

\[
\begin{cases}
\displaystyle y_{n+1}=y_n+hz_n+\frac{h^2}{6}(L_1+L_2+L_3),\\[6pt]
\displaystyle z_{n+1}=z_n+\frac h6(L_1+2L_2+2L_3+L_4),
\end{cases}
\]

这里

\[
L_1=f(x_n,y_n,z_n),
\]

\[
L_2=f\left(x_n+\frac h2,y_n+\frac h2z_n,z_n+\frac h2L_1\right),
\]

\[
L_3=f\left(x_n+\frac h2,y_n+\frac h2z_n+\frac{h^2}{4}L_1,z_n+\frac h2L_2\right),
\]

\[
L_4=f\left(x_n+h,y_n+hz_n+\frac{h^2}{2}L_2,z_n+hL_3\right).
\]

\end{document}
